\section{Executive Summary}
Payment fraud is a major risk for e-commerce platforms because even a small percentage of fraudulent transactions can lead to significant financial losses and damage customer trust. In this project, our team develops a real-time fraud detection system that helps an e-commerce platform identify suspicious transactions and decide whether each transaction should be approved, manually reviewed, or automatically blocked.

The project uses the required Kaggle Online Payments Fraud Detection Dataset and extends it with synthetic contextual data, such as customer account age, device/browser information, address mismatch indicators, failed payment attempts, and time-of-day behavior. This additional data is designed to make the fraud detection problem more realistic and closer to real business conditions.

From a business perspective, the main objective is not only to maximize model accuracy, but also to balance two competing goals: preventing fraudulent transactions and minimizing unnecessary friction for legitimate customers. Therefore, the project evaluates model performance using business-relevant KPIs, including fraud rate, precision, recall, false positive rate, financial loss avoided, cost of false alarms, and manual review rate.

The final output of the project is a fraud risk scoring system. Based on the predicted fraud probability, transactions are categorized into three risk levels. Low-risk transactions are approved automatically, medium-risk transactions are sent to manual review, and high-risk transactions are automatically blocked. This risk-based policy allows the platform to reduce financial loss while maintaining a reasonable customer experience.